\section*{अध्याय \ref{ch:g3:flowers-fruits-seeds} --- फूल, फल और बीज}

\begin{solution}{exo:g3:flowers-fruits-seeds:1}
\omterm{def:g3:flowers-fruits-seeds:parts}{पंखुड़ियाँ} (रंगीन और महकती हुई, वे \omterm{prop:g3:sorting-living-things:groups}{कीटों} को खींचती हैं); \omterm{def:g3:flowers-fruits-seeds:parts}{पुंकेसर} (पतली
डंडियाँ, जो \omterm{def:g3:flowers-fruits-seeds:parts}{पराग} लिए रहती हैं); और बीच में \omterm{def:g3:flowers-fruits-seeds:parts}{स्त्रीकेसर} (वह छोटी सुराही
जिसमें भावी \omterm{def:g2:from-seed-to-plant:seed}{बीज} रखे हैं)।
\end{solution}

\begin{solution}{exo:g3:flowers-fruits-seeds:2}
\omterm{def:g3:flowers-fruits-seeds:parts}{पराग} वह महीन पीला चूर्ण है जिसे \omterm{def:g3:flowers-fruits-seeds:parts}{पुंकेसर} लिए रहते हैं। परागण यानी \omterm{def:g3:flowers-fruits-seeds:parts}{पराग}
का \omterm{def:g3:flowers-fruits-seeds:parts}{स्त्रीकेसर} तक पहुँचना --- जिसे ज़्यादातर आने वाले \omterm{prop:g3:sorting-living-things:groups}{कीट} पहुँचाते हैं,
और कभी-कभी हवा।
\end{solution}

\begin{solution}{exo:g3:flowers-fruits-seeds:3}
(क) वे मुरझाकर गिर जाती हैं। (ख) वह फूलकर \omterm{prop:g3:flowers-fruits-seeds:transform}{फल} बन जाता है। (ग) वे \omterm{prop:g3:flowers-fruits-seeds:transform}{फल} के
भीतर \omterm{def:g2:from-seed-to-plant:seed}{बीज} बन जाते हैं।
\end{solution}

\begin{solution}{exo:g3:flowers-fruits-seeds:4}
अप्रैल: सफ़ेद \omterm{def:g3:flowers-fruits-seeds:parts}{फूल}, आती-जाती मधुमक्खियाँ --- परागण। मई: \omterm{def:g3:flowers-fruits-seeds:parts}{पंखुड़ियाँ}
गिरती हैं; हर परागित \omterm{def:g3:flowers-fruits-seeds:parts}{फूल} का \omterm{def:g3:flowers-fruits-seeds:parts}{स्त्रीकेसर} फूलकर नन्ही हरी गोली बन जाता
है। जून: वही गोली बढ़कर लाल, मीठी चेरी बन चुकी है, जिसके भीतर गुठली
--- यानी \omterm{def:g2:from-seed-to-plant:seed}{बीज} --- है।
\end{solution}

\begin{solution}{exo:g3:flowers-fruits-seeds:5}
क्योंकि \omterm{prop:g3:sorting-living-things:groups}{कीटों} के आए बिना ज़्यादातर फूलों का परागण होता ही नहीं, और बिना
परागण वाला \omterm{def:g3:flowers-fruits-seeds:parts}{फूल} \omterm{prop:g3:flowers-fruits-seeds:transform}{फल} बने बिना गिर जाता है। मधुमक्खियों की पराग-ढुलाई ही
फूलों को चेरियों में बदलती है।
\end{solution}

\begin{solution}{exo:g3:flowers-fruits-seeds:6}
उसे काटो और \omterm{def:g2:from-seed-to-plant:seed}{बीज} ढूँढ़ो: भीतर \omterm{def:g2:from-seed-to-plant:seed}{बीज} होने का मतलब है कि वह किसी \omterm{def:g3:flowers-fruits-seeds:parts}{स्त्रीकेसर}
से बना है --- यानी \omterm{prop:g3:flowers-fruits-seeds:transform}{फल}; \omterm{def:g2:from-seed-to-plant:seed}{बीज} न होने का मतलब है \omterm{def:g1:plants-around-us:plant}{पौधे} का कोई और भाग। टमाटर
\omterm{def:g2:from-seed-to-plant:seed}{बीजों} से भरा है: \omterm{prop:g3:flowers-fruits-seeds:transform}{फल}। गाजर में एक भी \omterm{def:g2:from-seed-to-plant:seed}{बीज} नहीं: \omterm{prop:g3:flowers-fruits-seeds:transform}{फल} नहीं, बल्कि \omterm{def:g1:plants-around-us:parts}{जड़}।
\end{solution}

\begin{solution}{exo:g3:flowers-fruits-seeds:7}
खीरा: \omterm{def:g2:from-seed-to-plant:seed}{बीज} --- \omterm{prop:g3:flowers-fruits-seeds:transform}{फल}। सलाद-पत्ता: कोई \omterm{def:g2:from-seed-to-plant:seed}{बीज} नहीं --- \omterm{def:g1:plants-around-us:parts}{पत्तियाँ}, \omterm{prop:g3:flowers-fruits-seeds:transform}{फल} नहीं। हरी
फली: क़तार में \omterm{def:g2:from-seed-to-plant:seed}{बीज} --- \omterm{prop:g3:flowers-fruits-seeds:transform}{फल}। आलू: कोई \omterm{def:g2:from-seed-to-plant:seed}{बीज} नहीं --- ज़मीन के नीचे फूला
हुआ \omterm{def:g1:plants-around-us:parts}{डंठल}, \omterm{prop:g3:flowers-fruits-seeds:transform}{फल} नहीं।
\end{solution}

\begin{solution}{exo:g3:flowers-fruits-seeds:8}
उत्तर खुला है। आम चौंकाने वाली बातें: टमाटर, खीरा और फली का \omterm{prop:g3:flowers-fruits-seeds:transform}{फल} निकलना;
और स्ट्रॉबेरी का अपने नन्हे \omterm{def:g2:from-seed-to-plant:seed}{बीज} बाहर की ओर लिए रहना।
\end{solution}

\begin{solution}{exo:g3:flowers-fruits-seeds:9}
आलूबुख़ारे लगभग नहीं होंगे। \omterm{prop:g3:flowers-fruits-seeds:transform}{फल} सिर्फ़ परागित \omterm{def:g3:flowers-fruits-seeds:parts}{फूल} के \omterm{def:g3:flowers-fruits-seeds:parts}{स्त्रीकेसर} से ही
आ सकता है; अप्रैल में \omterm{def:g3:flowers-fruits-seeds:parts}{फूल} मर जाने पर बदलने को कुछ बचा ही नहीं, गर्मी
चाहे कितनी भी अच्छी रहे।
\end{solution}

\begin{solution}{exo:g3:flowers-fruits-seeds:10}
\omterm{def:g1:plants-around-us:tree}{पेड़} के लिए सेब \omterm{def:g2:from-seed-to-plant:seed}{बीज} बनाने और उनकी रक्षा करने का एक तरीक़ा है। उसके
बीचोंबीच दाने पड़े हैं --- यानी \omterm{def:g2:from-seed-to-plant:seed}{बीज} --- और उनमें से हर दाना बोया जाए तो
अंकुरित होकर सेब का नया \omterm{def:g1:plants-around-us:tree}{पेड़} बन सकता है। पैकिंग हमें स्वादिष्ट लगती है,
यह हमारी क़िस्मत है (और \omterm{def:g1:plants-around-us:tree}{पेड़} के लिए अपने \omterm{def:g2:from-seed-to-plant:seed}{बीजों} को घर से दूर पहुँचाने का
एक ज़रिया)।
\end{solution}
